\section{Dependencies and quantifiers of the canonical construction}
\label{app:v108-audit}
Theorem~\ref{thm:v108-pool} uses only the inherited Borel physical
pairing space, finite one-sided endpoint traces, the regular-segment
variation bound, the one-copy uncut volume, and the fixed-window
paired and uncut average estimates. Its construction supplies the
comparison weights explicitly; overlap one is an identity, not a
new geometric assumption. Theorem~\ref{thm:v108-maximal} concerns
exactly Definition~\ref{def:v108-certificate}, at fixed $r,\Delta,K$.
It makes no claim of optimality over arbitrary comparison geometries,
transport costs, or changing windows. Theorem~\ref{thm:v108-dominance}
is a pointwise theorem, so it does not confuse optimization of total
angular mass with optimization in the required roof topology.

All source multipliers are evaluated at the original target roof.
At each fixed collision count, radius and exact labels, integration
and the positive physical sum precede every essential supremum.
The outer stage fixes $\varepsilon(B),h_B,r_B,\Delta_B,K_B$.
After those choices come the radius and exact-label suprema and the
collision limsup; $B\to\infty$ is last. Finite-count exhaustion in
Proposition~\ref{prop:v108-exhaustion} is not substituted for this
ordered estimate. No inverse-residence or logarithmic moment is
claimed for the Lorentz source by the new construction.

The original singular trace extensions still require their stated
physical hypotheses. An endpoint is not admitted to the comparison
space merely because it has been given a singular or stationary
name. Any source not paid by admissible capacity remains, with its
original positive weight, in \eqref{eq:v108-residual}. Thus the
construction addresses comparison selection and overlap, but the
collision-uniform shortage estimate \eqref{eq:v108-final-height}
requires a further proof on the original Lorentz family.
